\documentclass[10pt,a4paper]{amsart}
\usepackage[margin=19mm]{geometry}
\usepackage{amsmath,amssymb,mathtools}
\usepackage[scr=rsfs,cal=euler]{mathalfa}
\usepackage{xfrac}
\usepackage[unicode,hidelinks,bookmarks=false]{hyperref}
\newcommand{\ie}{i.e.\ }
\newcommand{\cf}{cf.\ }
\newcommand{\resp}{respectively\ }
\newcommand{\Subsection}{\subsection}
\providecommand{\nobreakdash}{\nobreak\hbox{-}}
\setlength{\emergencystretch}{2em}
\multlinegap0pt
\usepackage[shortlabels]{enumitem}
\usepackage{bm}
\usepackage{float}
\usepackage{amssymb}
\newtheorem{theorem}{Theorem}
\newtheorem{claim}{Claim}
\DeclareMathOperator*{\E}{\mathbb{E}}
\DeclareMathOperator{\HH}{H}
\DeclareMathOperator{\I}{I}
\DeclareMathOperator{\KL}{D_{\text{KL}}}
\newcommand{\MAXS}{\boldsymbol{MaxS}}
\newcommand{\MID}{\boldsymbol{MID}}
\newcommand{\PC}{\boldsymbol{PC}}
\newcommand{\POST}{\boldsymbol{post}}
\newcommand{\SECS}{{\mathcal{D}}}
\newcommand{\SUC}{\boldsymbol{suc}}
\DeclareMathOperator{\Supp}{\text{Supp}}
\newcommand{\TR}{\boldsymbol{tr}}
\title{Non-Adaptive Cryptanalytic Time-Space Lower Bounds via a Shearer-like Inequality for Permutations}
\author{Itai Dinur, Nathan Keller, Avichai Marmor}
\date{Source: arXiv:2505.00894v3 (CC BY 4.0)}
\begin{document}
\maketitle

\noindent\emph{Source and licence.} Non-Adaptive Cryptanalytic Time-Space Lower Bounds via a Shearer-like Inequality for Permutations. Version arXiv:2505.00894v3 (CC BY 4.0), distributed by arXiv under the Creative Commons Attribution 4.0 International licence, http://creativecommons.org/licenses/by/4.0/, as recorded in the arXiv licence metadata for this exact version and shown on the abstract page https://arxiv.org/abs/2505.00894v3. Full licence notice: \texttt{licenses/arxiv-2505.00894-CC-BY-4.0.txt}.\par\smallskip

\noindent\textit{Editorial scope.} Counted unit: the article's theorem middle (source0.tex lines 1208--1218). The article's complete original proof of it is delivered with it (source0.tex lines 1442--1532, one \texttt{proof} environment of 91 lines, which the article places away from the statement). No mathematics is written, repaired or re-derived: the statement and the proof are the article's own lines, in the article's own notation, and no retained line is substituted or re-flowed.  Retained uncounted as the source-local context the proof needs: lines 506--547; lines 1278--1284; lines 1333--1345.  Nothing was omitted from any retained span, so the package carries no omission record.  No retained line contains a citation, so the wrapper carries no bibliography and no bibliography entry.  No retained line contains a figure environment, an image-inclusion command or drawing code, so no image is delivered and the package declares no asset dependency.  Editorial formatting only, in addition: an AMS \texttt{amsart} preamble that restates the article's own declarations and macros used by the retained lines, namely the environments \texttt{theorem}, \texttt{claim} on separate counters, together with the standard packages those lines need.  The retained environments print in this document as Theorem 1, Claim 1 in order of appearance, whereas the article prints the same items with the numbering of its own full document.  The complete native source of the article as distributed in the arXiv e-print for this exact version is bundled unchanged as \texttt{sources/177457/source0.tex}.
\begin{enumerate}[label=(\arabic*), ref=\text{Property (}\arabic*\text{)}]
  \item \label{prope:conditioned entropy} Conditioning does not increase entropy, namely
$$\HH(X) \geq \HH(X \mid Y ),$$
   with equality if and only if $X,Y$ are independent, i.e., $P_{X,Y} = P_X \times P_Y$.
  %
  \item \label{prope:chain rule entropy} The \emph{chain rule for entropy} asserts that $$\HH(X,Y) = \HH(X) + \HH(Y \mid X).$$
Thus, $\HH(X,Y) = \HH(X) + \HH(Y)$ if and only if $X,Y$ are independent.
  % 
  \item \label{prope:entropy bound} The entropy of $X$ is upper-bounded by the logarithm of the size of its support, namely
  $$\HH(X) \leq \log|\Supp(\mathcal{X})|.$$
  %
  \item \label{prope:kl nonnegative} KL-divergence is non-negative, namely $$\KL(P_{X} \| Q_{X}) \geq 0.$$
  %
  \item \label{prope:kl convex} KL-divergence is convex, namely for $0 \leq \lambda \le 1$,
  $$\KL( \lambda P_X + (1 - \lambda) P'_X \| \lambda Q_X + (1 - \lambda)Q'_X ) \leq
  \lambda \KL( P_X \| Q_X ) + (1 - \lambda) \KL( P'_X \| Q'_X ).$$
  %
  \item \label{prope:kl and entropy} If $Q_X$ is uniform  over its support (which includes the support of $P_X$), then
    \begin{align*}
    \KL( P_X \| Q_X ) = &\; \E_{P_X(x)}[\log(1/ Q_X(x))] - \HH(P_X) = \E_{Q_X(x)}[\log(1/ Q_X(x))] - \HH(P_X) \\
    = &\; 
    \HH(Q_X) - \HH(P_X).
    \end{align*}
  %
  \item \label{prope:chain rule kl} The \emph{chain rule for KL-divergence} asserts that
  $$\KL( P_{X,Y} \| Q_{X,Y}) =  \KL( P_X \| Q_X ) + \KL( P_{Y \mid X} \| Q_{Y \mid X} ) .$$
  %
  \item \label{prope:data processing}
  The \emph{data processing inequality} asserts that for a function $f: \mathcal{X} \mapsto \mathcal{Y}$,
  $$\KL( P_X \| Q_X) \geq \KL( P_{f(X)} \| Q_{f(X)} ).$$
  %
\item \label{prope:pinsker}
A special case of \emph{Pinsker's inequality} states that
  $$2 (p - q)^2 \leq \KL(p \| q ).$$
  %
  \item \label{prope:information bound} The mutual information satisfies
  \begin{align*}
  &\; \I(X ; Y) = \KL(P_{X,Y} \| P_X P_Y) = 
  \KL(P_{X \mid Y} \| P_X) = \HH(X) - \HH(X \mid Y) \leq 
  \HH(X).
  \end{align*}
\end{enumerate}

\begin{claim} We have
\label{clm:basic with advice}
$$
q' 
\leq 
\widehat{\MAXS}(T).$$
\end{claim}

\begin{claim} \label{clm:advice main-eq2}
For any $z,\sigma^{in}$ in the support of $P_{Z,\Sigma^{in}}$,
\begin{align*} 
p_{z, \sigma^{in} }
\leq 
\min \left(
2 q'_{z, \sigma^{in} }
+ 
\frac{4 \kappa'_{z, \sigma^{in} } T_2 }{u},
q'_{z,\sigma^{in} } + \sqrt{\frac{\kappa'_{z,  \sigma^{in} } T_2} {{u}}} 
\right).
\end{align*}
\end{claim}

\begin{theorem} \label{thm:pc to middle}
Let $\PC := \PC(N,\SECS,\mathcal{M},\TR,\POST,\SUC)$ be a permutation challenge game, such that 
$\TR$ is $u$-uniform.
Let $(A_0,A_1)$ be an $(S,T)$ non-adaptive algorithm with preprocessing for $\PC$.
Denote by $\widehat{\MAXS}(T)$ the optimal success probability (with respect to $\SUC$) of an adversary to $\MID$ with running time $T$.
Then, the success probability of $(A_0,A_1)$ is at most
\[
\min \left(2\cdot \widehat{\MAXS}(T) + \frac{4\log(2) S T }{u}, \widehat{\MAXS}(T) + 
\sqrt{\frac{\log(2) ST} {{u}}}\right).
\]
\end{theorem}

\begin{proof} [of Theorem~\ref{thm:pc to middle}]
We average both sides of the inequality of Claim~\ref{clm:advice main-eq2} 
over $P_{Z,\Sigma^{in} }(z,\sigma^{in}) = Q'_{Z,\Sigma^{in} }(z,\sigma^{in})$.

On the left-hand-side, we obtain
$
\E_{P_{Z,\Sigma^{in}}(z, \sigma^{in} )}
[p_{z,\sigma^{in} }]
= p$.

On the right-hand-side,
we consider $q'_{z, \sigma^{in}}$ and $\kappa'_{z,\sigma^{in} }$ separately.

First, since
$P_{Z,\Sigma^{in} }(z,\sigma^{in}) = Q'_{Z, \Sigma^{in} }(z,\sigma^{in})$, we have
$$
\E_{Q'_{Z,\Sigma^{in}}(z,\sigma^{in})}[
q'_{z, \sigma^{in}} ] =
q' \leq 
\widehat{\MAXS}(T),$$
where the inequality is by Claim~\ref{clm:basic with advice}.

Second, we have
\begin{align*} 
\E_{Q_{Z,\Sigma^{in}}(z, \sigma^{in} )} [ \kappa'_{z,\sigma^{in} } ] 
&=  
\E_{Q_{Z,\Sigma^{in}}(z, \sigma^{in} )} [ \KL( P_{ \Sigma_{\overline{\mathcal{S}}} \mid \Sigma^{in} = \sigma^{in} , Z = z} \| Q'_{ \Sigma_{ \overline{\mathcal{S}} } \mid \Sigma^{in} = \sigma^{in} } ) ]. 
\end{align*}
Recall that 
for any $\sigma^{in}$
$,
Q'_{ \Sigma_{ \overline{\mathcal{S}} } \mid \Sigma^{in} = \sigma^{in} }$
is uniform 
over its support, hence, by~\ref{prope:kl and entropy}
\begin{align*} 
& \E_{Q_{Z,\Sigma^{in}}(z, \sigma^{in} )} [ \KL( P_{ \Sigma_{\overline{\mathcal{S}}} \mid \Sigma^{in} = \sigma^{in} , Z = z } \| Q'_{ \Sigma_{ \overline{\mathcal{S}} } \mid \Sigma^{in} = \sigma^{in} } ) ] \\
&= 
\E_{Q_{Z,\Sigma^{in}}(z, \sigma^{in} )}[\HH(Q'_{\Sigma_{ \overline{\mathcal{S}} } \mid \Sigma^{in} = \sigma^{in}} ) - 
\HH(P_{ \Sigma_{\overline{\mathcal{S}}} \mid \Sigma^{in} = \sigma^{in} , Z = z } ) ] \\
&= 
\HH(Q'_{\Sigma_{ \overline{\mathcal{S}} } \mid \Sigma^{in}} ) - 
\HH(P_{ \Sigma_{\overline{\mathcal{S}}} \mid \Sigma^{in} , Z} )
= 
\HH(Q'_{\Sigma_{ \overline{\mathcal{S}} } \mid \mathcal{S},\Sigma_{\mathcal{S} },Z } ) - 
\HH(P_{ \Sigma_{\overline{\mathcal{S}}} \mid \mathcal{S},\Sigma_{\mathcal{S} },Z} ) \\
&=
\HH(Q'_{\Sigma_{ \overline{\mathcal{S}} } \mid \Sigma_{\mathcal{S}},Z } ) - 
\HH(P_{ \Sigma_{\overline{\mathcal{S}}} \mid \Sigma_{\mathcal{S}},Z} ).
\end{align*} 
By~\ref{prope:chain rule entropy} and~\ref{prope:conditioned entropy},
\begin{align*} 
\HH(Q'_{\Sigma_{ \overline{\mathcal{S}} } \mid \Sigma_{\mathcal{S}},Z } ) 
&= 
\HH(Q'_{ \Sigma \mid Z} ) - 
\HH(Q'_{ \Sigma_{\mathcal{S}} \mid Z}) = 
\HH(Q'_{ \Sigma \mid Z} ) - 
\HH(P_{ \Sigma_{\mathcal{S}} \mid Z}) \leq
\HH(Q'_{ \Sigma} ) - 
\HH(P_{ \Sigma_{\mathcal{S}} \mid Z}) \\
&\leq 
\HH(P_{ \Sigma} ) - 
\HH(P_{ \Sigma_{\mathcal{S}} \mid Z}),
\end{align*} 

and similarly,
$$
\HH(P_{ \Sigma_{\overline{\mathcal{S}}} \mid \Sigma_{\mathcal{S}},Z} ) = 
\HH(P_{ \Sigma \mid Z} ) - 
\HH(P_{ \Sigma_{\mathcal{S}} \mid Z}). 
$$
Therefore, by~\ref{prope:information bound},~\ref{prope:entropy bound},
and since $(A_0,A_1)$ is an $(S,T)$ algorithm,
\begin{align*} 
\E_{Q_{Z,\Sigma^{in}}(z, \sigma^{in} )} [ \kappa'_{z,\sigma^{in} } ] \leq
\HH(P_{\Sigma}) - \HH(P_{\Sigma \mid Z} ) = \I(\Sigma ; Z) \leq
\HH(Z) \leq 
\log|\Supp(\mathcal{Z})| \leq \log(2) S.
\end{align*}

Overall, we obtain 
$$ 
p \leq 
2 \cdot \widehat{\MAXS}(T) + \frac{\log(2) 4 S T }{u},
\textup{ and }
p \leq 
\widehat{\MAXS}(T) + 
\sqrt{\frac{\log(2) ST} {{u}}},
$$
where the second inequality uses the concavity of the square root function.
This concludes the proof.
\end{proof}

\end{document}
